\documentclass[11pt,a4paper]{article}
\usepackage[T1]{fontenc}
\usepackage[utf8]{inputenc}
\usepackage{lmodern,amsmath,amssymb}
\usepackage[margin=25mm]{geometry}
\usepackage[hidelinks]{hyperref}
\hypersetup{pdftitle={The SU(2) bridge midpoint in closed form: the curvature term is exact},pdfauthor={navstokgap research working notes}}
\newcommand{\tightlist}{\setlength{\itemsep}{0pt}\setlength{\parskip}{0pt}}
\setlength{\emergencystretch}{3em}
\setcounter{secnumdepth}{0}
\begin{document}
\hypertarget{the-su2-bridge-midpoint-in-closed-form-the-curvature-term-is-exact}{%
\section{The SU(2) bridge midpoint in closed form: the curvature term is
exact}\label{the-su2-bridge-midpoint-in-closed-form-the-curvature-term-is-exact}}

\textbf{Result, 2026-09-27.} Proposition 6 of the
\href{series-parallel-gauge-refinement.md}{series/parallel note}
described the midpoint of the Brownian bridge on \(SU(2)\), the building
block of the non-abelian parallel insertion, only to leading order in
the heat time, with the curvature factor \(1/h\), \(h=(d/4)\cot(d/4)\),
derived by Laplace's method and labelled as asymptotics. For the
fundamental character the midpoint expectation is computable exactly.
For a bridge of duration \(t\) from \(e\) to a group element at geodesic
distance \(\theta\in[0,2\pi)\),

\[E\,\chi_{1/2}(m)=e^{-t/32}\Bigl[2\cos\frac\theta4-\frac t{2\theta}\,
\sin\frac\theta4\Bigr]\ \ \text{up to image terms of relative size}\
\Bigl(2+\frac{32\pi^2}t\Bigr)e^{-4\pi(2\pi-\theta)/t}\]

(Theorem 1). Expanded to first order it is exactly the prediction of
Proposition 6(b),
\(2\cos\frac\theta4\,[1-\frac t{32}(1+\frac2{h(\theta)})]\), because
\(\frac t{32}\cdot\frac2h=\frac t{4\theta}\tan\frac\theta4\). So the
curvature term N2, the side-flux-dependent softening of the mid-face
weight, is confirmed with no remainder of order \(t^2\) in the \(w=0\)
sector: the only correction to the displayed formula is the
exponentially small winding of the bridge around the group. This is the
first exact non-abelian statement about the parallel insertion in three
dimensions, and it settles the fundamental-representation part of the
uniform Laplace remainder that Proposition 7 of the series/parallel note
assumes.

The ingredients are character orthogonality and Poisson summation; the
computation is elementary and no novelty is claimed for it.

\hypertarget{setting}{%
\subsection{1. Setting}\label{setting}}

Group metric as in the series/parallel note §1: \(C_2(j)=j(j+1)\),
\(SU(2)\) the three-sphere of radius 2, heat kernel
\(k_t=\sum_j(2j+1)e^{-tj(j+1)/2}\chi_j\) with
\(\chi_j(\theta)=\sin((2j+1)\theta/2)/\sin(\theta/2)\) at rotation angle
\(\theta\) (the geodesic distance from \(e\)). The midpoint \(m\) of the
bridge of duration \(t\) from \(e\) to \(g\) has density
\(k_{t/2}(m)\,k_{t/2}(m^{-1}g)/k_t(g)\) with respect to Haar measure.
Its geodesic midpoint \(m_*=\exp(\frac12\log g)\) has rotation angle
\(\theta/2\), so
\(\chi_{1/2}(m_*)=\sin(\theta/2)/\sin(\theta/4)=2\cos(\theta/4)\).

\hypertarget{the-exact-formula}{%
\subsection{2. The exact formula}\label{the-exact-formula}}

\textbf{Theorem 1.} Put
\(\Theta_\pm(\theta)=\sum_{w\in\mathbb Z}(\pm1)^w e^{-(\theta-4\pi w)^2/(2t)}\)
and
\(\Xi_\pm(\theta)=\sum_{w\in\mathbb Z}(\pm1)^w(\theta-4\pi w) e^{-(\theta-4\pi w)^2/(2t)}\).
For \(\theta\in(0,2\pi)\),

\[E\,\chi_{1/2}(m)=e^{-t/32}\,
\frac{2\cos(\theta/4)\,\Xi_-(\theta)-\frac t2\sin(\theta/4)\,\Theta_-(\theta)}
{\Xi_+(\theta)}. \tag{1}\]

Keeping only \(w=0\) gives the displayed formula of the summary.

\emph{Proof.} Expand both heat kernels in characters. Character
orthogonality in the form
\(\int\chi_a(m)\chi_b(m^{-1}g)\,dm=\delta_{ab}\chi_a(g)/d_a\), together
with \(\chi_{j_1}\chi_{1/2}=\chi_{j_1+1/2}+\chi_{j_1-1/2}\), gives

\[k_t(g)\,E\,\chi_{1/2}(m)=\sum_{j_1}\ \sum_{j_2=j_1\pm1/2}(2j_1+1)\,
e^{-t[C_2(j_1)+C_2(j_2)]/4}\,\chi_{j_2}(g).\]

Write \(n=2j_1+1\). Then
\(C_2(j_1)+C_2(j_2)=\frac14[n^2+(n\pm1)^2-2] =\frac12(n\pm\frac12)^2-\frac38\),
so the exponential is \(e^{3t/32}e^{-t(n\pm1/2)^2/8}\). Multiply by
\(\sin(\theta/2)\), reindex the \(j_2=j_1-\frac12\) sum by
\(n\mapsto n+1\), and write \(s=n+\frac12\) over
\(s\in\frac12+\mathbb N\): the summand becomes
\((s-\frac12)\sin\frac{(s+\frac12)\theta}2+(s+\frac12)\sin\frac{(s-\frac12)\theta}2 =2s\cos\frac\theta4\sin\frac{s\theta}2-\sin\frac\theta4\cos\frac{s\theta}2\),
which is even in \(s\). Hence

\[k_t(g)\sin\frac\theta2\;E\,\chi_{1/2}(m)=\frac{e^{3t/32}}2
\Bigl[2\cos\frac\theta4\,A-\sin\frac\theta4\,B\Bigr],\]

with \(B=\sum_{s\in\frac12+\mathbb Z}e^{-ts^2/8}\cos\frac{s\theta}2\)
and
\(A=\sum_{s\in\frac12+\mathbb Z}s\,e^{-ts^2/8}\sin\frac{s\theta}2=-2\partial_\theta B\).
In the same way \(k_t(g)\sin\frac\theta2=\frac{e^{t/8}}2A_0\) with
\(A_0=\sum_{n\in\mathbb Z}n\,e^{-tn^2/8}\sin\frac{n\theta}2\). Poisson
summation over the half-integers (which introduces the sign \((-1)^w\))
and over the integers gives \(B=\sqrt{8\pi/t}\,\Theta_-\),
\(A=(2/t)\sqrt{8\pi/t}\,\Xi_-\) and \(A_0=(2/t)\sqrt{8\pi/t}\,\Xi_+\).
Divide. \(\square\)

\textbf{The image terms.} For \(\theta\in(0,2\pi)\) the pair \(w=\pm1\)
in \(\Xi_\pm\) has, relative to the \(w=0\) term
\(\theta e^{-\theta^2/(2t)}\), the size
\(e^{-8\pi^2/t}\,|2\theta\cosh(4\pi\theta/t)\mp8\pi\sinh(4\pi\theta/t)|/\theta \le(2+32\pi^2/t)\,e^{-4\pi(2\pi-\theta)/t}\),
using \(\sinh x\le x\cosh x\). In \(\Theta_-\) the same pair has
relative size at most \(2e^{-4\pi(2\pi-\theta)/t}\). The pairs
\(|w|\ge2\) are smaller by a further factor \(e^{-16\pi^2/t}\) up to
polynomial factors. This is the large-field term N4, the winding of the
bridge around the group, and it is negligible uniformly on
\(\theta\le\theta_0<2\pi\).

\hypertarget{comparison-with-proposition-6}{%
\subsection{3. Comparison with Proposition
6}\label{comparison-with-proposition-6}}

Proposition 6(b) predicts that \(m=m_*e^\xi\) with \(\xi\) Gaussian of
mean zero and covariance \(\frac t4{\rm diag}(1,\frac1h,\frac1h)\) in
the frame (axis, transverse). In the spin-\(\frac12\) representation
\((\xi\cdot T)^2=-\frac14|\xi|^2\), so
\(E\,D^{1/2}(e^\xi)=1-\frac18E|\xi|^2+O(t^2)=1-\frac t{32}(1+\frac2h)+O(t^2)\),
and
\(E\,\chi_{1/2}(m)=2\cos\frac\theta4\,[1-\frac t{32}(1+\frac2h)]+O(t^2)\).
The \(w=0\) part of (1) equals
\(2\cos\frac\theta4\,e^{-t/32}[1-\frac t{4\theta}\tan\frac\theta4]\),
and
\(\frac t{32}\cdot\frac2{h(\theta)}=\frac t{16}\cdot\frac{\tan(\theta/4)}{\theta/4} =\frac t{4\theta}\tan\frac\theta4\).
The two agree at first order, curvature factor included. As
\(\theta\to0\) both give \(2(1-\frac{3t}{32})\), the flat value
\(E|\xi|^2=\frac{3t}4\).

Formula (1) says more than the asymptotics. In the \(w=0\) sector the
fundamental character of the midpoint is the product of an isotropic
factor \(e^{-t/32}\) and the exactly linear curvature correction
\(1-\frac t{4\theta}\tan\frac\theta4\); there are no higher-order terms
in \(t\). The correction grows as \(\theta\to2\pi\), where
\(\tan(\theta/4)\) diverges and the midpoint becomes bimodal, and there
the image terms take over. This is the non-abelian counterpart of the
parity factor \(\rho_t(\phi)\) of the \(U(1)\) cube (Proposition 3 of
the series/parallel note).

\hypertarget{consequence-for-state}{%
\subsection{4. Consequence for STATE}\label{consequence-for-state}}

The curvature term of the \(SU(2)\) parallel insertion is now exact for
the fundamental character. Two steps remain before Proposition 7 becomes
a theorem for the isolated \(SU(2)\) cube: the same computation for
every spin \(J\), since the cube's character sum
\(\sum_Jd_Je^{-t_mC_2(J)/2}{\rm tr}_J\prod_jE\,D^J(m_j^{\pm1})\)
involves all spins, and the matrix structure of \(E\,D^J\) beyond its
trace, which Proposition 6(a) reduces to one real diagonal matrix per
cut face. The first follows the proof above with
\(\chi_{j_1}\chi_J=\sum_{|j_1-J|}^{j_1+J}\chi_{j_2}\) and is the next
step on this line.
\end{document}
